\chapter{Implémentation}

\section{Environnement technique}

\subsection{Infrastructure de calcul}

L'ensemble des expériences a été réalisé sur la plateforme Kaggle Notebooks, qui fournit un environnement de calcul reproductible avec les caractéristiques suivantes :

\begin{table}[H]
\centering
\caption{Spécifications de l'environnement d'exécution}
\label{tab:env}
\begin{tabular}{ll}
\toprule
\textbf{Composant} & \textbf{Spécification} \\
\midrule
Plateforme        & Kaggle Notebooks \\
GPU               & NVIDIA Tesla T4 $\times$ 2 (16 GB VRAM chacun) \\
RAM               & 29 GB (mode GPU) \\
Stockage          & 20 GB (espace de travail) \\
CPU               & 4 cœurs Intel Xeon \\
\bottomrule
\end{tabular}
\end{table}

\subsection{Stack technologique}

\begin{table}[H]
\centering
\caption{Stack logicielle du projet}
\label{tab:stack}
\begin{tabular}{ll}
\toprule
\textbf{Composant} & \textbf{Technologie / Version} \\
\midrule
Langage              & Python 3.10 \\
Traitement distribué & Apache Spark 3.5 (PySpark) \\
Deep Learning        & TensorFlow 2.15 / Keras \\
Machine Learning     & XGBoost 2.0, scikit-learn 1.3 \\
Visualisation        & Matplotlib 3.7, Seaborn 0.12, Plotly 5.15 \\
Dashboard            & Streamlit 1.28+ \\
Format de données    & Parquet (PySpark), CSV (pandas) \\
Stockage modèles     & .keras (DNN), .json (XGBoost), .pkl (scaler) \\
\bottomrule
\end{tabular}
\end{table}

\subsection{Structure du projet}

\begin{lstlisting}[language=bash, caption=Structure du projet, label=lst:structure, style=pythonstyle]
pfe/
|-- notebook/               # Notebooks Kaggle (Phase 1-6)
|   |-- 01_setup_and_data_loading.ipynb
|   |-- 02_preprocessing_and_eda.ipynb
|   |-- 03_feature_engineering.ipynb
|   |-- 04_dnn_training.ipynb
|   |-- 05_xgboost_spark_training.ipynb
|   +-- 06_evaluation_comparison.ipynb
|-- streamlit_dashboard/    # Dashboard interactif
|   |-- app.py              # Page d'accueil
|   |-- pages/
|   |   |-- 1_EDA.py
|   |   |-- 2_Model_Results.py
|   |   +-- 3_Spark_Impact.py
|   |-- utils/helpers.py    # Chargement des artifacts
|   |-- results/            # JSON + CSV + PNG
|   +-- data/               # brfss_sample.csv
|-- phases/                 # Documentation phases
|-- checkpoints/            # Rapports d'execution
+-- rapport/                # Ce rapport LaTeX
    |-- main.tex
    |-- references.bib
    |-- chapters/
    +-- figures/
\end{lstlisting}

\section{Phase 1 --- Chargement et nettoyage des données (PySpark)}

\subsection{Initialisation de la session Spark}

\begin{lstlisting}[caption=Initialisation de Spark, label=lst:spark_init]
from pyspark.sql import SparkSession
import pyspark.sql.functions as F

spark = SparkSession.builder \
    .appName("BRFSS_CVD_Prediction") \
    .master("local[*]") \
    .config("spark.driver.memory", "8g") \
    .config("spark.sql.shuffle.partitions", "200") \
    .getOrCreate()

# Chargement du dataset pooled ML
df_raw = spark.read.parquet(
    "/kaggle/input/brfss-2020-2024-dataset/"
    "brfss_2020_2024_pooled_ml.parquet"
)
print(f"Dimensions brutes : {df_raw.count():,} lignes x "
      f"{len(df_raw.columns)} colonnes")
# Output : 2,176,776 lignes x 129 colonnes
\end{lstlisting}

\subsection{Pipeline de binarisation BRFSS}

Le format BRFSS encode les réponses binaires avec 1=Oui et 2=Non, contrairement à la convention 0/1 attendue par les modèles de ML. Une étape de transformation explicite est nécessaire avant toute opération.

\begin{lstlisting}[caption=Binarisation et création de la variable cible, label=lst:binarize]
# Variable cible : _MICHD (coronary heart disease OR myocardial infarction)
df = df_raw.withColumn('HeartDiseaseorAttack',
    F.when(F.col('_MICHD') == 1, 1.0)
     .when(F.col('_MICHD') == 2, 0.0)
     .otherwise(None))

# Variables binaires : 1->1, 2->0
for col_name in ['CVDSTRK3', 'DIFWALK', 'EXERANY2']:
    new_name = {'CVDSTRK3': 'Stroke', 'DIFWALK': 'DiffWalk',
                'EXERANY2': 'PhysActivity'}[col_name]
    df = df.withColumn(new_name,
        F.when(F.col(col_name) == 1, 1.0)
         .when(F.col(col_name) == 2, 0.0)
         .otherwise(None))

# BMI : division par 100
df = df.withColumn('BMI', F.col('_BMI5') / 100.0)

# Age : tranche 14 (ne sait pas) -> null
df = df.withColumn('Age',
    F.when(F.col('_AGEG5YR') == 14, None)
     .otherwise(F.col('_AGEG5YR').cast('float')))
\end{lstlisting}

\subsection{Déduplication et statistiques}

\begin{lstlisting}[caption=Déduplication et diagnostic, label=lst:dedup]
SELECTED_COLS = ['HeartDiseaseorAttack', 'BMI', 'Smoker', 'Stroke',
                 'Diabetes', 'PhysActivity', 'GenHlth', 'MentHlth',
                 'PhysHlth', 'DiffWalk', 'Sex', 'Age',
                 'Education', 'HvyAlcoholConsump']

df_selected = df.select(SELECTED_COLS).dropDuplicates()
n_after = df_selected.count()  # -> 1,760,241

# Distribution de la cible
df_selected.groupBy('HeartDiseaseorAttack').count().show()
# +---------------------+--------+
# |HeartDiseaseorAttack |count   |
# +---------------------+--------+
# |0.0                  |1578452 |
# |1.0                  |181789  |
# +---------------------+--------+
# Prevalence : 181789 / 1760241 = 10.32%
\end{lstlisting}

\section{Phase 2 --- Feature Engineering (pandas + scikit-learn)}

\subsection{Conversion et imputation}

Pour les phases de feature engineering et d'entraînement des modèles, le DataFrame Spark est converti en DataFrame pandas après filtrage et nettoyage :

\begin{lstlisting}[caption=Imputation des valeurs manquantes, label=lst:imputation]
import pandas as pd
import numpy as np

df_pd = df_selected.toPandas()

# Imputation : médiane pour les continus, mode pour les catégoriels
for col in ['BMI', 'MentHlth', 'PhysHlth']:
    median_val = df_pd[col].median()
    df_pd[col].fillna(median_val, inplace=True)

for col in ['Smoker', 'Stroke', 'Diabetes', 'PhysActivity',
            'DiffWalk', 'Sex', 'GenHlth', 'Age', 'Education',
            'HvyAlcoholConsump']:
    mode_val = df_pd[col].mode()[0]
    df_pd[col].fillna(mode_val, inplace=True)
\end{lstlisting}

\subsection{Création des features dérivées}

\begin{lstlisting}[caption=Features dérivées RiskScore et HealthIndex, label=lst:derived]
# RiskScore : somme des facteurs de risque binaires disponibles
df_pd['RiskScore'] = (df_pd['Smoker'] + df_pd['Diabetes'] +
                      df_pd['Stroke'] + df_pd['DiffWalk'] +
                      df_pd['HvyAlcoholConsump'])

# HealthIndex : etat de sante global normalise
df_pd['HealthIndex'] = (df_pd['GenHlth'] +
                        df_pd['MentHlth'] / 6.0 +
                        df_pd['PhysHlth'] / 6.0)
\end{lstlisting}

\section{Phase 3 --- Entraînement du DNN (TensorFlow/Keras)}

\subsection{Préparation des données}

\begin{lstlisting}[caption=Préparation du jeu de données pour le DNN, label=lst:dnn_prep]
from sklearn.model_selection import train_test_split
from sklearn.preprocessing import StandardScaler
from sklearn.utils.class_weight import compute_class_weight

FEATURES = ['Age', 'GenHlth', 'Stroke', 'DiffWalk', 'Diabetes',
            'PhysHlth', 'Smoker', 'Sex', 'PhysActivity',
            'HvyAlcoholConsump', 'BMI', 'Education', 'MentHlth',
            'RiskScore', 'HealthIndex']

X = df_pd[FEATURES].values
y = df_pd['HeartDiseaseorAttack'].values

X_train, X_test, y_train, y_test = train_test_split(
    X, y, test_size=0.30, random_state=42, stratify=y)

# Normalisation
scaler = StandardScaler()
X_train = scaler.fit_transform(X_train)
X_test  = scaler.transform(X_test)

# Class weights (rapport ~9:1)
class_weights = compute_class_weight('balanced',
    classes=np.unique(y_train), y=y_train)
cw_dict = {0: class_weights[0], 1: class_weights[1]}
# -> {0: 0.558, 1: 4.847}
\end{lstlisting}

\subsection{Architecture et entraînement}

\begin{lstlisting}[caption=Construction et entraînement du DNN, label=lst:dnn_train]
import tensorflow as tf
from tensorflow.keras import layers, callbacks

def build_dnn(input_dim):
    inp = tf.keras.Input(shape=(input_dim,))
    x = layers.Dense(256, activation='relu')(inp)
    x = layers.BatchNormalization()(x)
    x = layers.Dropout(0.3)(x)
    x = layers.Dense(128, activation='relu')(x)
    x = layers.BatchNormalization()(x)
    x = layers.Dropout(0.3)(x)
    x = layers.Dense(64, activation='relu')(x)
    out = layers.Dense(1, activation='sigmoid')(x)
    return tf.keras.Model(inp, out)

model = build_dnn(len(FEATURES))
model.compile(
    optimizer=tf.keras.optimizers.Adam(learning_rate=1e-3),
    loss='binary_crossentropy',
    metrics=['accuracy', tf.keras.metrics.AUC(name='auc')]
)

cbs = [
    callbacks.EarlyStopping(monitor='val_auc', patience=5,
                             restore_best_weights=True, mode='max'),
    callbacks.ReduceLROnPlateau(monitor='val_auc', factor=0.5,
                                patience=3, mode='max'),
    callbacks.ModelCheckpoint('dnn_best_model.h5',
                               monitor='val_auc', save_best_only=True)
]

history = model.fit(
    X_train, y_train,
    epochs=50, batch_size=1024,
    validation_split=0.20,
    class_weight=cw_dict,
    callbacks=cbs, verbose=1
)
# Résultat : 36 epochs (arret a l'epoch 36, meilleur a l'epoch 31)
# Best val_auc = 0.8205
\end{lstlisting}

\subsection{Optimisation du seuil de classification}

\begin{lstlisting}[caption=Optimisation du seuil via courbe Précision-Rappel, label=lst:threshold]
from sklearn.metrics import precision_recall_curve

y_pred_proba = model.predict(X_test).flatten()

# Courbe precision-rappel
precisions, recalls, thresholds = precision_recall_curve(
    y_test, y_pred_proba)

# F1 par seuil
f1_scores = (2 * precisions * recalls /
             (precisions + recalls + 1e-8))

best_idx   = np.argmax(f1_scores[:-1])
best_thresh = float(thresholds[best_idx])
# -> seuil optimal = 0.6713

y_pred_optimal = (y_pred_proba >= best_thresh).astype(int)
\end{lstlisting}

\section{Phase 4 --- XGBoost et Benchmark Spark}

\subsection{Optimisation des hyperparamètres XGBoost}

\begin{lstlisting}[caption=RandomizedSearchCV pour XGBoost, label=lst:xgb_search]
from xgboost import XGBClassifier
from sklearn.model_selection import RandomizedSearchCV

param_dist = {
    'n_estimators':    [200, 300, 400, 500],
    'max_depth':       [3, 4, 5, 6, 7],
    'learning_rate':   [0.01, 0.05, 0.1],
    'subsample':       [0.7, 0.8, 0.9, 1.0],
    'colsample_bytree':[0.6, 0.7, 0.8, 0.9],
    'min_child_weight':[1, 3, 5, 7],
    'gamma':           [0, 0.1, 0.2, 0.5],
}

xgb_base = XGBClassifier(
    scale_pos_weight=len(y_train[y_train==0])/len(y_train[y_train==1]),
    eval_metric='auc', use_label_encoder=False,
    tree_method='gpu_hist', device='cuda'
)

rs = RandomizedSearchCV(
    xgb_base, param_dist,
    n_iter=20, cv=3, scoring='roc_auc',
    n_jobs=-1, random_state=42, verbose=2
)
rs.fit(X_train, y_train)
# Meilleure AUC CV : 0.8211
# Temps de recherche : 887 secondes
\end{lstlisting}

\subsection{Benchmark Spark vs. XGBoost}

\begin{lstlisting}[caption=Benchmark de performance par fraction, label=lst:benchmark]
import time
from pyspark.ml.classification import GBTClassifier as SparkGBT
from pyspark.ml.feature import VectorAssembler

fractions = [0.10, 0.25, 0.50, 0.75, 1.00]
results = []

for frac in fractions:
    df_frac = df_train_spark.sample(fraction=frac, seed=42)
    n_rows = df_frac.count()

    # --- Spark GBT ---
    assembler = VectorAssembler(inputCols=FEATURES,
                                outputCol='features')
    df_vec = assembler.transform(df_frac)
    gbt_spark = SparkGBT(maxIter=100, maxDepth=6,
                         featuresCol='features',
                         labelCol='HeartDiseaseorAttack')
    t0 = time.time()
    model_spark = gbt_spark.fit(df_vec)
    spark_time  = time.time() - t0

    # --- XGBoost standalone ---
    X_frac = X_train[:int(frac * len(X_train))]
    y_frac = y_train[:int(frac * len(y_train))]
    xgb_bench = XGBClassifier(**rs.best_params_,
                              tree_method='hist')
    t0 = time.time()
    xgb_bench.fit(X_frac, y_frac)
    standalone_time = time.time() - t0

    results.append({'fraction': frac, 'n_rows': n_rows,
                    'spark_time': spark_time,
                    'standalone_time': standalone_time})
\end{lstlisting}

\section{Phase 6 --- Dashboard Streamlit}

Le tableau de bord est une application Streamlit statique qui charge les artefacts pré-calculés (JSON, CSV, PNG) depuis le dossier \texttt{results/}. Il ne nécessite aucune inférence ML ni aucun accès à Spark. L'architecture est conçue pour fonctionner sur n'importe quelle machine disposant de Python 3.10+.

\begin{lstlisting}[caption=Lancement du dashboard, language=bash, style=pythonstyle]
# Installation
pip install -r streamlit_dashboard/requirements.txt

# Lancement
cd streamlit_dashboard
python -m streamlit run app.py --server.port 8501
\end{lstlisting}

\textbf{Structure du dashboard :}
\begin{itemize}
    \item \textbf{Page d'accueil} : 4 KPIs (observations, variables, meilleur AUC, prévalence), description du projet et navigation.
    \item \textbf{Page EDA} : distribution de la cible, distributions univariées interactives, heatmap de corrélation Pearson, boxplots des outliers, taux de MCV par sous-groupe. Filtres par sexe et tranche d'âge.
    \item \textbf{Page Résultats Modèles} : KPIs comparatifs, tableau des métriques, bar chart, courbes ROC superposées, matrices de confusion, radar chart, importance XGBoost, courbes d'apprentissage DNN, analyse du seuil, critères PRD.
    \item \textbf{Page Impact Spark} : KPIs de temps d'entraînement, benchmark scalabilité (échelle log), avantage XGBoost (Nx plus rapide), comparaison AUC, conclusions.
\end{itemize}
